Táto práca skúma, či je možné vytvoriť praktického Linux \gls{cli} asistenta
pomocou malých lokálnych \glspl{llm} pri offline a hardvérových obmedzeniach.
Práca navrhuje shAI, modulárny prototyp založený na \gls{rag}, ktorý pozostáva z
klasifikátora typu príkaz-vs.-otázka, vyhľadávača dokumentov indexujúceho man
stránky Linuxu a sumarizátora, ktorý generuje stručné odpovede zo získaných
pasáží. Implementácia využíva Ollama na lokálnu inferenciu a \gls{faiss} na
vektorové indexovanie.

Klasifikátor bol testovaný na vyváženej vzorke dát \texttt{classifier-100};
presnosť sa pohybovala od 50 \% do 88 \%, pričom najlepší model vyžadoval 2,9
sekundy na jeden dopyt. Presnosť vyhľadávača bola vyhodnocovaná na kolekciách s
10 a 100 dokumentmi (\texttt{sample-10} a \texttt{sample-100}) s využitím
rôznych prístupov k vyhľadávaniu. Prístup na úrovni dokumentov dosiahol
úspešnosť v top-5 nájdených dokumentoch od 86,3 \% do 92,0 \% s hodnotou MRR od
0,724 do 0,806 na kolekcii so 100 dokumentmi, pričom však
generoval obrovské výstupné veľkosti (až do 120-tisíc tokenov). Na vyriešenie
tohto problému boli implementované prístupy vyhľadávania na úrovni dokumentov s
preusporiadaním na úrovni pasáží a samostatné vyhľadávanie na
úrovni pasáží. Tieto prístupy vykazovali oveľa kompaktnejšie veľkosti
vygenerovaného výstupu (~5-tisíc tokenov) a dosiahli potrebnú presnosť.

Latencia aplikácie bola vyhodnocovaná na datasetoch \texttt{sample-10},
\texttt{sample-100} a \texttt{sample-1000}. Ukázalo sa, že vyhľadávanie na
úrovni pasáží a navrhované vyhľadávanie na úrovni dokumentov s preusporiadaním
na úrovni pasáží dosiahli najlepší celkový výkon v závislosti od veľkosti
zvoleného embedding modelu, s latenciou pod 2 sekundy pre vyhľadávanie a pod 10
sekúnd pre celý proces RAG.

Výsledky ukazujú, že architektúry postavené primárne na vyhľadávaní
(retrieval-first) sú pre lokálnu pomoc v prostredí CLI praktické, zatiaľ čg
otvorené agentové prístupy zostávajú pri malých modeloch
nespoľahlivé. Najlepšou voľbou pre techniku vyhľadávania je vyhľadávanie
na úrovni pasáží. Práca poskytuje open-source prototyp a empirické referenčné
test , ktoré slúžia ako návod pre budúci výskum v oblasti dizajnu
úloh s obmedzenými zdrojmi.
